\chapter{图网络、表示学习与智能推荐系统}

\label{chap:12}

到第 \ref{chap:11} 章为止，本书已经把无穷维凸优化的主要技术链搭建完成：第 \ref{chap:03} 章给出 Hilbert 空间与谱论工作台，第 \ref{chap:05} 章和第 \ref{chap:06} 章给出次微分与共轭对偶语言，第 \ref{chap:09} 章和第 \ref{chap:10} 章把这些对象组织成可计算的算子分裂与 ADMM 模板，第 \ref{chap:11} 章则说明它们如何落到控制与金融建模。本章继续沿着这条主线向前推进，不过对象从控制轨道与财富过程换成了现代学习系统中的函数表示、图结构与分布式一致性。

本章沿四条线索展开。第 \ref{sec:12-1} 节讨论 RKHS 与 Representer Theorem，说明为什么“在无限维函数空间学习”仍可坍缩为有限样本系数问题。第 \ref{sec:12-2} 节把图拉普拉斯、Dirichlet 能量与消息传递放进谱论和邻近算子的语言，解释图学习中真正可被凸优化准确接口化的部分。第 \ref{sec:12-3} 节转向多塔表示与匹配问题，用核化或 lifted 的凸 surrogate 解释双塔检索、对比学习与影子价格。最后，第 \ref{sec:12-4} 节把第 \ref{chap:10} 章的 ADMM 与原--对偶分裂落到推荐系统中的共识训练、分片嵌入与图正则模型。

本章不是独立的机器学习教材，也不会把一般深度非凸训练硬写成凸优化。我们只处理那些能够被无穷维凸分析、谱论、单调算子或 ADMM 语言稳定刻画的结构：核方法、图平滑、凸匹配 surrogate 和分布式一致性。Boyd 的有限维表达依旧是重要母语，但在这里，它们被用来说明欧氏模型如何坍缩自更一般的函数空间与算子框架，而不是作为唯一舞台。

\section{再生核 Hilbert 空间}

\label{sec:12-1}

本节的任务是把表示学习中最稳定的一条主线写回泛函分析：不是直接讨论网络参数，而是先讨论函数类本身。RKHS 的重要性在于，它允许我们一方面保留无限维函数空间的表达力，另一方面又通过样本点评价把问题压缩到有限参数。这个压缩并不是经验技巧，而是 Hilbert 空间几何、直接法存在性和正则化共同作用的结果。

\subsection{正定核与 RKHS}

\begin{definition}[正定核]\label{def:positive-definite-kernel}
设 $\mathcal X$ 为非空集合。函数
\[
K\colon \mathcal X\times \mathcal X\to \mathbb R
\]
称为 $\mathcal X$ 上的\emph{正定核}，如果对任意有限点集 $x_1,\dots,x_m\in \mathcal X$ 与任意系数 $c_1,\dots,c_m\in \mathbb R$，都有
\[
\sum_{i,j=1}^m c_i c_j K(x_i,x_j)\ge 0.
\]
\end{definition}

\begin{definition}[再生核 Hilbert 空间]\label{def:rkhs}
设 $K$ 为定义~\ref{def:positive-definite-kernel} 的正定核。若 Hilbert 空间 $\mathcal H_K$ 由定义在 $\mathcal X$ 上的实值函数组成，并满足：
\begin{enumerate}
  \item 对任意 $x\in \mathcal X$，函数 $K_x:=K(x,\cdot)$ 属于 $\mathcal H_K$；
  \item 对任意 $f\in \mathcal H_K$ 与任意 $x\in \mathcal X$，
  \[
  f(x)=\langle f,K_x\rangle_{\mathcal H_K},
  \]
\end{enumerate}
则称 $\mathcal H_K$ 为核 $K$ 的\emph{再生核 Hilbert 空间}。
\end{definition}

\begin{proposition}[点评价泛函连续]\label{prop:rkhs-evaluation-continuity}
在定义~\ref{def:rkhs} 的条件下，对任意 $x\in \mathcal X$，点评价映射
\[
\delta_x\colon \mathcal H_K\to \mathbb R,\qquad \delta_x(f)=f(x)
\]
是连续线性泛函，并满足估计
\[
|f(x)|\le \|f\|_{\mathcal H_K}\sqrt{K(x,x)},
\qquad \forall f\in \mathcal H_K.
\]
\end{proposition}

\begin{proof}
线性性显然。由再生性质，
\[
f(x)=\langle f,K_x\rangle_{\mathcal H_K}.
\]
再由 Cauchy--Schwarz 不等式，
\[
|f(x)|\le \|f\|_{\mathcal H_K}\|K_x\|_{\mathcal H_K}.
\]
而
\[
\|K_x\|_{\mathcal H_K}^2
=
\langle K_x,K_x\rangle_{\mathcal H_K}
=
K(x,x),
\]
故结论成立。
\end{proof}

\begin{remark}
命题~\ref{prop:rkhs-evaluation-continuity} 把 RKHS 与第 \ref{sec:3-1} 节定义~\ref{def:hilbert-space} 及定理~\ref{thm:riesz-representation-hilbert} 直接接了起来：样本点评价之所以能够被写成内积，不是因为数据科学中特殊的坐标技巧，而是因为在 Hilbert 空间里连续线性泛函本来就能由向量表示。
\end{remark}

\subsection{Tikhonov 正则化与表示定理}

在样本 $x_1,\dots,x_m\in \mathcal X$ 上，考虑带 Tikhonov 项的学习问题
\[
\min_{f\in \mathcal H_K}
\left\{
\Phi\bigl(f(x_1),\dots,f(x_m)\bigr)+\lambda \|f\|_{\mathcal H_K}^2
\right\},
\qquad \lambda>0,
\]
其中 $\Phi\colon \mathbb R^m\to \mathbb R\cup\{+\infty\}$ 为 proper、凸且下半连续泛函。由于正则项强制了 $\mathcal H_K$ 范数，结合命题~\ref{prop:rkhs-evaluation-continuity} 的连续性，上式正是第 \ref{sec:5-1} 节定理~\ref{thm:direct-method-reflexive} 可处理的 Hilbert 空间直接法模型。

\begin{theorem}[Representer Theorem]\label{thm:representer-theorem}
设 $\Phi$ 与 $\lambda$ 如上。若 $f_\star$ 是问题
\[
\min_{f\in \mathcal H_K}
\left\{
\Phi\bigl(f(x_1),\dots,f(x_m)\bigr)+\lambda \|f\|_{\mathcal H_K}^2
\right\}
\]
的极小元，则存在系数 $\alpha_1,\dots,\alpha_m\in \mathbb R$ 使
\[
f_\star(\cdot)=\sum_{i=1}^m \alpha_i K(x_i,\cdot).
\]
换言之，极小元一定落在有限维子空间
\[
\operatorname{span}\{K(x_1,\cdot),\dots,K(x_m,\cdot)\}
\]
中。
\end{theorem}

\begin{proof}
记
\[
M:=\operatorname{span}\{K(x_1,\cdot),\dots,K(x_m,\cdot)\}\subset \mathcal H_K.
\]
由第 \ref{sec:3-1} 节定理~\ref{thm:orthogonal-projection}，任意 $f\in \mathcal H_K$ 都可唯一分解为
\[
f=f_M+f_\perp,
\qquad
f_M\in M,\quad f_\perp\in M^\perp.
\]
由于 $f_\perp\perp K(x_i,\cdot)$，对每个样本点 $x_i$ 有
\[
f_\perp(x_i)=\langle f_\perp,K(x_i,\cdot)\rangle_{\mathcal H_K}=0.
\]
因此
\[
f(x_i)=f_M(x_i),\qquad i=1,\dots,m,
\]
从而数据项只依赖于 $f_M$：
\[
\Phi\bigl(f(x_1),\dots,f(x_m)\bigr)
=
\Phi\bigl(f_M(x_1),\dots,f_M(x_m)\bigr).
\]
另一方面，由正交分解，
\[
\|f\|_{\mathcal H_K}^2=\|f_M\|_{\mathcal H_K}^2+\|f_\perp\|_{\mathcal H_K}^2
\ge \|f_M\|_{\mathcal H_K}^2.
\]
故对任意 $f\in \mathcal H_K$，
\[
\Phi\bigl(f(x_1),\dots,f(x_m)\bigr)+\lambda \|f\|_{\mathcal H_K}^2
\ge
\Phi\bigl(f_M(x_1),\dots,f_M(x_m)\bigr)+\lambda \|f_M\|_{\mathcal H_K}^2.
\]
这说明把 $f$ 替换成其在 $M$ 上的正交投影不会增大目标值。于是任何极小元都可取在 $M$ 中。由于 $M$ 由核截面张成，结论成立。
\end{proof}

\begin{remark}
定理~\ref{thm:representer-theorem} 是本章最关键的有限样本坍缩结论。它告诉我们：无限维函数学习并没有被偷偷地变成有限维，而是在样本驱动与二次正则的共同作用下，自然塌缩到了一个由数据决定的有限维截面空间。若把目标写成
\[
F(f):=\Phi(Cf)+\lambda\|f\|_{\mathcal H_K}^2,
\]
则最优性条件可以进一步写成
\[
0\in C^\ast\partial \Phi(Cf_\star)+2\lambda f_\star,
\]
这正是第 \ref{sec:5-3} 节定义~\ref{def:subdifferential} 的函数空间版本。
\end{remark}

\subsection{典型例子}

\begin{example}[核岭回归]\label{ex:kernel-ridge-regression}
取
\[
\Phi(z_1,\dots,z_m)=\frac{1}{m}\sum_{i=1}^m (z_i-y_i)^2.
\]
则相应问题为
\[
\min_{f\in \mathcal H_K}
\left\{
\frac{1}{m}\sum_{i=1}^m (f(x_i)-y_i)^2+\lambda \|f\|_{\mathcal H_K}^2
\right\}.
\]
由定理~\ref{thm:representer-theorem}，极小元可写成
\[
f_\star(\cdot)=\sum_{i=1}^m \alpha_i K(x_i,\cdot).
\]
把它代回目标，即得到关于系数向量 $\alpha$ 的有限维二次规划。把这个例子放在这里，是为了说明 RKHS 并不排斥可计算性；相反，它给出了核岭回归这类算法的最自然理论母体。
\end{example}

\begin{example}[有限维线性回归的坍缩]
若核取为线性核
\[
K(x,x')=\langle x,x'\rangle_{\mathbb R^d},
\]
则 $\mathcal H_K$ 与 $\mathbb R^d$ 等同，函数 $f$ 由向量 $w$ 表示为 $f(x)=\langle w,x\rangle$。此时定理~\ref{thm:representer-theorem} 退化为
\[
w_\star\in \operatorname{span}\{x_1,\dots,x_m\},
\]
这正是有限维线性回归的经典结论。
\end{example}

\begin{example}[文本与政策匹配中的核表示]
在政策文本检索或企业画像匹配中，常把文档先映射到高维语义特征，再以核函数近似相似性。若学习目标是正则化经验风险，则定理~\ref{thm:representer-theorem} 说明最优匹配函数只由训练样本对应的核截面张成。把这个例子放在这里，是为了给第 \ref{sec:12-3} 节的多塔匹配问题准备一个函数空间入口。
\end{example}

\subsection*{有限维对照}

Boyd 在有限维正则化模型里常把岭回归写成矩阵方程或二次规划；这当然是正确的，但它默认参数已经是一个向量。本节的结论说明，Boyd 的有限维公式并不是与函数学习平行的一套理论，而是 RKHS 正则化经过定理~\ref{thm:representer-theorem} 坍缩后的欧氏版本。正因为如此，后文谈图表示、检索匹配与分布式训练时，才能不断把现代学习对象重新翻译回 Hilbert 空间与凸优化语法。

\subsection*{习题（待补）}

\begin{exercise}[表示定理占位]
补一题：把定理~\ref{thm:representer-theorem} 中的二次正则替换成严格单调的范数函数，说明表示定理仍如何成立。
\end{exercise}

\begin{exercise}[核岭回归系数占位]
补一题：对例~\ref{ex:kernel-ridge-regression} 写出 Gram 矩阵形式的系数方程，并说明它如何退化为有限维岭回归。
\end{exercise}


\section{图神经网络的泛函算子解释}

\label{sec:12-2}

图神经网络最常见的工程描述是“聚合邻居再更新表示”，但若只停留在这个层面，很多关键结构都会被遮蔽：为什么图卷积总要与拉普拉斯或归一化邻接联系在一起，为什么平滑、扩散和低通滤波会反复出现，以及哪些图学习步骤真的能够被第 \ref{chap:09} 章和第 \ref{chap:10} 章的算子语言准确接口化。本节选择从图拉普拉斯与 Dirichlet 能量出发，抓住那些可被谱论、prox 和分裂方法稳定描述的部分。

\subsection{图信号、拉普拉斯与能量}

\begin{definition}[图拉普拉斯算子]\label{def:graph-laplacian-operator}
设 $\mathscr G=(V,E,w)$ 为有限加权无向图，其中 $V=\{1,\dots,n\}$，边权满足 $w_{ij}=w_{ji}\ge 0$。记图信号空间
\[
H_V:=\mathbb R^n
\]
并赋予标准内积。定义图拉普拉斯算子 $L_{\mathscr G}\colon H_V\to H_V$ 为
\[
(L_{\mathscr G}u)_i:=\sum_{j=1}^n w_{ij}(u_i-u_j),
\qquad i=1,\dots,n.
\]
\end{definition}

\begin{definition}[图 Dirichlet 能量]\label{def:graph-dirichlet-energy}
在定义~\ref{def:graph-laplacian-operator} 的条件下，称泛函
\[
\mathcal E_{\mathscr G}(u)
:=
\frac12\sum_{i,j=1}^n w_{ij}(u_i-u_j)^2
\]
为图 $\mathscr G$ 上信号 $u$ 的\emph{Dirichlet 能量}。
\end{definition}

\begin{proposition}[图拉普拉斯的正性]\label{prop:graph-laplacian-positive}
算子 $L_{\mathscr G}$ 是 $H_V$ 上的自伴正半定算子，并满足
\[
\langle L_{\mathscr G}u,u\rangle=\mathcal E_{\mathscr G}(u),
\qquad \forall u\in H_V.
\]
\end{proposition}

\begin{proof}
由权重对称性，
\[
\langle L_{\mathscr G}u,v\rangle
=
\sum_{i=1}^n \sum_{j=1}^n w_{ij}(u_i-u_j)v_i
=
\frac12\sum_{i,j=1}^n w_{ij}(u_i-u_j)(v_i-v_j),
\]
交换 $u,v$ 可知它自伴。令 $v=u$，得到
\[
\langle L_{\mathscr G}u,u\rangle
=
\frac12\sum_{i,j=1}^n w_{ij}(u_i-u_j)^2
=
\mathcal E_{\mathscr G}(u)\ge 0.
\]
故 $L_{\mathscr G}$ 正半定。
\end{proof}

\begin{remark}
由于 $H_V$ 是有限维 Hilbert 空间，$L_{\mathscr G}$ 同时还是紧自伴算子。于是第 \ref{sec:3-4} 节定理~\ref{thm:compact-selfadjoint-spectral} 可直接应用于它：图傅里叶基、特征值平滑强度和谱滤波，本质上都是紧自伴谱理论在图上的有限维坍缩。
\end{remark}

\subsection{消息传递与图滤波}

\begin{proposition}[消息传递是图滤波的一种线性原型]\label{prop:message-passing-graph-filter}
设 $p$ 为一元多项式，则算子
\[
T_p:=p(L_{\mathscr G})
\]
与 $L_{\mathscr G}$ 共享同一组特征向量，因此是一个图谱滤波器。特别地，线性消息传递层
\[
u^{+}=(I-\tau L_{\mathscr G})u,
\qquad 0<\tau<\|L_{\mathscr G}\|^{-1},
\]
对应于一阶低通图滤波；若再附加逐点非线性与通道权矩阵，就得到常见 GNN 层的线性骨架。
\end{proposition}

\begin{proof}
由定理~\ref{thm:compact-selfadjoint-spectral}，存在正交基 $\{e_k\}_{k=1}^n$ 与实特征值 $\lambda_k\ge 0$，使
\[
L_{\mathscr G}e_k=\lambda_k e_k.
\]
于是
\[
T_p e_k=p(\lambda_k)e_k,
\]
故 $T_p$ 与 $L_{\mathscr G}$ 同时对角化。特别地，当 $p(\lambda)=1-\tau\lambda$ 时，
\[
T_p=I-\tau L_{\mathscr G},
\]
它衰减高频特征模态，正是一阶低通滤波。
\end{proof}

\begin{theorem}[图扩散与学习接口]\label{thm:graph-diffusion-learning-interface}
设 $\tau>0$，$y\in H_V$。则问题
\[
\min_{u\in H_V}
\left\{
\frac12\|u-y\|^2+\tau \mathcal E_{\mathscr G}(u)
\right\}
\]
有唯一极小元 $u_\tau$，并满足
\[
u_\tau=(I+\tau L_{\mathscr G})^{-1}y=\operatorname{prox}_{\tau\mathcal E_{\mathscr G}}(y).
\]
此外，映射
\[
y\longmapsto (I-\tau L_{\mathscr G})y
\]
是上述 resolvent 的一阶显式近似，因此线性图扩散、图平滑和一类 GNN 消息传递层都可被看成第 \ref{sec:10-1} 节定义~\ref{def:forward-backward-iteration} 在图信号空间上的具体接口。
\end{theorem}

\begin{proof}
由命题~\ref{prop:graph-laplacian-positive}，
\[
\mathcal E_{\mathscr G}(u)=\langle L_{\mathscr G}u,u\rangle
\]
是凸二次型。故目标是严格凸连续函数，唯一极小元存在。对极小元 $u_\tau$ 写一阶最优性条件，
\[
0=u_\tau-y+\tau L_{\mathscr G}u_\tau.
\]
于是
\[
(I+\tau L_{\mathscr G})u_\tau=y,
\qquad
u_\tau=(I+\tau L_{\mathscr G})^{-1}y.
\]
另一方面，由定义~\ref{def:proximal-mapping}，$\operatorname{prox}_{\tau\mathcal E_{\mathscr G}}(y)$ 正是上式的唯一极小元，故
\[
u_\tau=\operatorname{prox}_{\tau\mathcal E_{\mathscr G}}(y).
\]
最后，由 resolvent 展开
\[
(I+\tau L_{\mathscr G})^{-1}y
=
y-\tau L_{\mathscr G}y+O(\tau^2),
\]
可见显式一步 $y\mapsto y-\tau L_{\mathscr G}y$ 是对应的前向近似。
\end{proof}

\begin{remark}
定理~\ref{thm:graph-diffusion-learning-interface} 只 formalize 那些可嵌入凸能量或算子分裂语言的图学习部分。一般深层 GNN 训练包含参数共享、非线性组合与非凸优化；本节并不把这些内容伪装成凸问题，而是只抽取其中最稳定的图滤波、平滑与扩散骨架。
\end{remark}

\subsection{典型例子}

\begin{example}[图平滑最小化]
若 $y$ 是带噪观测信号，则定理~\ref{thm:graph-diffusion-learning-interface} 给出的极小元平衡了数据保真和图平滑。这正是半监督学习、标签传播与图正则回归中的基础模型。
\end{example}

\begin{example}[邻接矩阵与归一化拉普拉斯的有限维坍缩]
在有限图上，常见 GCN 层写作
\[
H^{+}=\sigma(\widetilde D^{-1/2}\widetilde A\widetilde D^{-1/2}HW).
\]
把
\[
\widetilde L:=I-\widetilde D^{-1/2}\widetilde A\widetilde D^{-1/2}
\]
代入，就得到
\[
\widetilde D^{-1/2}\widetilde A\widetilde D^{-1/2}=I-\widetilde L,
\]
因此消息传递矩阵本质上就是一阶图滤波器。
\end{example}

\begin{example}[知识图谱与 item-item 推荐中的关系传播]
在知识图谱推荐或 item-item 推荐中，边权编码相似度、共现或关系强度。若只保留图传播的线性骨架，则每一步更新都可解释为对邻域一致性和局部平滑的近似求解。把这个例子放在这里，是为了把 GNN 直接接到推荐系统的图结构接口上。
\end{example}

\subsection*{有限维对照}

Boyd 不会直接讲 GNN，但他关于二次型、谱分解和分布式平滑的有限维母语，在这里仍然有效。区别只在于：工程上常把“邻接矩阵乘一次”当成一个架构模块，而本节说明它其实是图拉普拉斯谱滤波、Dirichlet 能量最小化以及 prox/前向步的一种坍缩实现。理解这一点以后，第 \ref{sec:12-4} 节的分布式图推荐模型就不再只是经验配方，而能重新接回 ADMM 与原--对偶分裂。

\subsection*{习题（待补）}

\begin{exercise}[图拉普拉斯正性占位]
补一题：证明命题~\ref{prop:graph-laplacian-positive}，并说明常数信号为什么总落在零特征值对应空间中。
\end{exercise}

\begin{exercise}[图扩散与 prox 占位]
补一题：从定理~\ref{thm:graph-diffusion-learning-interface} 出发，说明图扩散为何可以被写成定义~\ref{def:proximal-mapping} 的邻近算子。
\end{exercise}


\section{多塔模型与无限维特征匹配}

\label{sec:12-3}

推荐与检索系统中的双塔、 多塔模型通常以工程术语出现：用户塔、物品塔、内积打分、负样本、对比学习。若直接沿着这条叙事线展开，很容易把关键数学结构埋进网络细节里。本节的做法是先把匹配问题写成 Hilbert 空间中的配对风险，再说明当我们采用核化特征、固定编码器或 lifted convex surrogate 时，这些工程对象如何进入第 \ref{chap:06} 章的对偶和第 \ref{chap:07} 章的 KKT 语言。

\subsection{双塔表示与配对风险}

\begin{definition}[双塔嵌入模型]\label{def:two-tower-embedding-model}
设 $\mathcal X,\mathcal Y$ 为两类对象的样本空间，$\mathcal H_X,\mathcal H_Y$ 为对应的 RKHS，核分别为 $K_X,K_Y$。记张量积 Hilbert 空间
\[
\mathcal H:=\mathcal H_X\otimes \mathcal H_Y.
\]
对任意 $M\in \mathcal H$，定义匹配得分
\[
s_M(x,y)
:=
\left\langle
M,\,
K_X(x,\cdot)\otimes K_Y(y,\cdot)
\right\rangle_{\mathcal H}.
\]
若
\[
M=\sum_{r=1}^R u_r\otimes v_r,
\]
则
\[
s_M(x,y)=\sum_{r=1}^R u_r(x)v_r(y),
\]
这给出一个 $R$ 维双塔或多塔型表示。
\end{definition}

\begin{definition}[配对损失泛函]\label{def:pairwise-loss-functional}
给定训练样本
\[
\{(x_i,y_i,\varepsilon_i)\}_{i=1}^m\subset \mathcal X\times \mathcal Y\times \{-1,1\},
\]
以及凸损失函数 $\ell_i\colon \mathbb R\to \mathbb R\cup\{+\infty\}$，定义配对损失
\[
\mathcal L(M)
:=
\sum_{i=1}^m \ell_i\bigl(\varepsilon_i s_M(x_i,y_i)\bigr).
\]
其中 $\varepsilon_i=1$ 表示正匹配，$\varepsilon_i=-1$ 表示负匹配。
\end{definition}

\begin{definition}[对比风险泛函]\label{def:contrastive-risk-functional}
在定义~\ref{def:pairwise-loss-functional} 的背景下，给定正则参数 $\lambda>0$，称
\[
\mathcal R_\lambda(M)
:=
\mathcal L(M)+\frac{\lambda}{2}\|M\|_{\mathcal H}^2
\]
为\emph{对比风险泛函}。它把配对损失与 Tikhonov 正则统一到同一个 Hilbert 空间优化模板中。
\end{definition}

\subsection{表示定理与对偶解释}

\begin{proposition}[双塔模型的表示定理]\label{prop:two-tower-representer}
设 $\mathcal R_\lambda$ 如定义~\ref{def:contrastive-risk-functional}，并设其在 $\mathcal H$ 上有极小元 $M_\star$。则存在系数 $\alpha_1,\dots,\alpha_m\in \mathbb R$ 使
\[
M_\star
=
\sum_{i=1}^m \alpha_i
\bigl(K_X(x_i,\cdot)\otimes K_Y(y_i,\cdot)\bigr).
\]
也就是说，最优匹配算子总落在训练样本所张成的有限维张量子空间中。
\end{proposition}

\begin{proof}
张量积空间 $\mathcal H$ 仍是 Hilbert 空间。令
\[
\psi_i:=K_X(x_i,\cdot)\otimes K_Y(y_i,\cdot),
\qquad
M_0:=\operatorname{span}\{\psi_1,\dots,\psi_m\}.
\]
由定义~\ref{def:two-tower-embedding-model}，
\[
s_M(x_i,y_i)=\langle M,\psi_i\rangle_{\mathcal H}.
\]
因此对任意分解
\[
M=M_0^\parallel+M_0^\perp,
\qquad
M_0^\perp\in M_0^\perp,
\]
都有
\[
s_M(x_i,y_i)=s_{M_0^\parallel}(x_i,y_i),\qquad i=1,\dots,m.
\]
于是配对损失只依赖 $M_0^\parallel$。另一方面，
\[
\|M\|_{\mathcal H}^2
=
\|M_0^\parallel\|_{\mathcal H}^2+\|M_0^\perp\|_{\mathcal H}^2
\ge
\|M_0^\parallel\|_{\mathcal H}^2.
\]
故把 $M$ 替换为 $M_0^\parallel$ 不会增大定义~\ref{def:contrastive-risk-functional} 的值。于是极小元可以取在 $M_0$ 中。这个结论也可被视为定理~\ref{thm:representer-theorem} 在张量积 RKHS 上的直接应用。
\end{proof}

\begin{theorem}[正则化匹配模型的对偶表示]\label{thm:embedding-duality-regularized}
设 $C\colon \mathcal H\to \mathbb R^m$ 为采样算子
\[
CM:=\bigl(\varepsilon_1 s_M(x_1,y_1),\dots,\varepsilon_m s_M(x_m,y_m)\bigr),
\]
并设 $\Phi\colon \mathbb R^m\to \mathbb R\cup\{+\infty\}$ 为 proper、凸且下半连续泛函。考虑原问题
\[
\inf_{M\in \mathcal H}
\left\{
\Phi(CM)+\frac{\lambda}{2}\|M\|_{\mathcal H}^2
\right\}.
\]
则其对偶问题为
\[
\sup_{\alpha\in \mathbb R^m}
\left\{
-\Phi^\ast(\alpha)-\frac{1}{2\lambda}\|C^\ast \alpha\|_{\mathcal H}^2
\right\}.
\]
若资格条件满足，则原--对偶最优解 $(M_\star,\alpha_\star)$ 满足
\[
M_\star=-\frac{1}{\lambda}C^\ast\alpha_\star,
\qquad
\alpha_\star\in \partial \Phi(CM_\star).
\]
\end{theorem}

\begin{proof}
设
\[
f(M):=\frac{\lambda}{2}\|M\|_{\mathcal H}^2,
\qquad
g(z):=\Phi(z),
\qquad
A:=C.
\]
则原问题正是
\[
\inf_{M\in \mathcal H}\{f(M)+g(AM)\}.
\]
由于 $\mathcal H$ 是 Hilbert 空间，$f$ 的共轭为
\[
f^\ast(\Xi)=\frac{1}{2\lambda}\|\Xi\|_{\mathcal H}^2.
\]
于是由第 \ref{sec:6-4} 节定理~\ref{thm:fenchel-rockafellar-duality}，
\[
\inf_{M}\{f(M)+g(AM)\}
=
\sup_{\alpha\in\mathbb R^m}
\{-f^\ast(-A^\ast \alpha)-g^\ast(\alpha)\}.
\]
代入 $f^\ast$ 与 $g^\ast=\Phi^\ast$，得到
\[
\sup_{\alpha\in \mathbb R^m}
\left\{
-\Phi^\ast(\alpha)-\frac{1}{2\lambda}\|C^\ast\alpha\|_{\mathcal H}^2
\right\}.
\]
再由推论~\ref{cor:kkt-via-fenchel-subdifferential}，原--对偶最优性等价于
\[
-C^\ast\alpha_\star\in \partial f(M_\star),
\qquad
\alpha_\star\in \partial \Phi(CM_\star).
\]
而
\[
\partial f(M)=\{\lambda M\},
\]
故
\[
M_\star=-\frac{1}{\lambda}C^\ast\alpha_\star.
\]
结论成立。
\end{proof}

\begin{remark}
定理~\ref{thm:embedding-duality-regularized} 说明，对比学习中的负样本权重、困难样本放大和约束影子价格，都可以通过对偶变量 $\alpha$ 统一理解。这里的关键不是把所有深度 encoder 证明成凸的，而是识别出当特征固定、核化近似或 lifted surrogate 成立时，匹配问题的确能进入 Fenchel 与 KKT 框架。
\end{remark}

\subsection{典型例子}

\begin{example}[矩阵分解与双线性匹配的有限维坍缩]
若
\[
\mathcal H_X=\mathbb R^p,\qquad \mathcal H_Y=\mathbb R^q,
\]
则 $\mathcal H_X\otimes \mathcal H_Y$ 可识别为 $p\times q$ 矩阵空间，$M$ 退化为双线性打分矩阵，
\[
s_M(x,y)=x^\top My.
\]
工业界常见的双塔参数化 $M=UV^\top$ 是它的低秩非凸特化；而本节讨论的是对完整矩阵或其凸 surrogate 的优化。
\end{example}

\begin{example}[双塔检索的核化解释]
在用户检索和召回场景中，常把查询与文档分别映射到公共向量空间，再用内积评分。若我们固定编码特征，只优化最终匹配算子 $M$，则定义~\ref{def:contrastive-risk-functional} 就给出了该模型的凸风险版本。
\end{example}

\begin{example}[政策--企业匹配]
把政策文本与企业画像分别看成 $\mathcal X$ 与 $\mathcal Y$ 中的对象，匹配目标是学得一个高分表示“政策适配”的双塔得分函数。把这个例子放在这里，是为了强调第 \ref{sec:12-1} 节的 RKHS 入口并不是抽象装饰，而可以直接进入政企推荐和政策匹配等实际建模。
\end{example}

\subsection*{有限维对照}

Boyd 并不直接讲双塔检索，但他关于正则化经验风险、线性算子复合和对偶变量的有限维语言，在本节仍然完全适用。差别只在于：现代推荐系统把“参数向量”换成了表示函数、张量特征与配对算子。理解这一点以后，双塔、多塔、对比学习就不再只是模型工程术语，而是可以被看成 RKHS 表示定理加上 Fenchel 对偶后的有限维坍缩。

\subsection*{习题（待补）}

\begin{exercise}[张量表示定理占位]
补一题：证明命题~\ref{prop:two-tower-representer}，并写出最优算子所在的有限维张量子空间。
\end{exercise}

\begin{exercise}[匹配对偶占位]
补一题：对平方损失或 logistic 型对比学习 surrogate，写出定理~\ref{thm:embedding-duality-regularized} 的对偶问题，并解释对偶变量的负样本权重含义。
\end{exercise}


\section{大型系统的分布式求解}

\label{sec:12-4}

前两节把表示函数和图结构都放回了 Hilbert 空间与谱论语言；真正的大型推荐系统还需要最后一步：把这些模型拆到多机、多分片和多数据块上求解。此时关键对象不再只是单个目标函数，而是局部子问题、全局一致性约束和乘子更新。本节把这部分内容统一写成第 \ref{sec:10-3} 节 ADMM 的应用层接口。

\subsection{共识模型与分片嵌入}

\begin{definition}[分布式共识模型]\label{def:distributed-consensus-model}
设 $Z$ 为 Hilbert 空间，$Z_1,\dots,Z_N$ 为局部副本空间，$P_i\colon Z\to Z_i$ 为有界线性算子。若优化问题写成
\[
\min_{z\in Z,\ z_i\in Z_i}
\left\{
\sum_{i=1}^N F_i(z_i)+G(z)
:\ z_i=P_i z,\ i=1,\dots,N
\right\},
\]
其中 $F_i,G$ 为 proper、凸且下半连续泛函，则称其为一个\emph{分布式共识模型}。
\end{definition}

\begin{definition}[分片嵌入问题]\label{def:sharded-embedding-problem}
在定义~\ref{def:distributed-consensus-model} 的背景下，若 $z$ 表示全局 embedding 参数，$z_i$ 表示第 $i$ 个工作节点上的局部参数副本，而 $F_i$ 编码该节点的局部样本损失、图正则或匹配风险，则称该模型为\emph{分片嵌入问题}。
\end{definition}

\subsection{ADMM 形式与局部更新}

\begin{proposition}[共识模型的 ADMM 分裂]\label{prop:consensus-admm-splitting}
定义~\ref{def:distributed-consensus-model} 的问题可写成第 \ref{sec:7-2} 节定义~\ref{def:abstract-cone-program} 的线性约束标准型。对给定罚参数 $\rho>0$，其 ADMM 迭代可写为
\[
z_i^{k+1}\in
\arg\min_{z_i}
\left\{
F_i(z_i)+\langle y_i^k,z_i-P_i z^k\rangle+\frac{\rho}{2}\|z_i-P_i z^k\|^2
\right\},
\]
\[
z^{k+1}\in
\arg\min_z
\left\{
G(z)-\sum_{i=1}^N \langle y_i^k,P_i z\rangle
+\frac{\rho}{2}\sum_{i=1}^N \|z_i^{k+1}-P_i z\|^2
\right\},
\]
\[
y_i^{k+1}=y_i^k+\rho(z_i^{k+1}-P_i z^{k+1}),
\qquad i=1,\dots,N.
\]
\end{proposition}

\begin{proof}
令
\[
X:=Z_1\times \cdots\times Z_N,
\qquad
x:=(z_1,\dots,z_N),
\qquad
f(x):=\sum_{i=1}^N F_i(z_i).
\]
再令
\[
Bz:=(P_1 z,\dots,P_N z)\in X,
\qquad
g(z):=G(z).
\]
则原问题等价于
\[
\min_{x,z}\{f(x)+g(z):x-Bz=0\}.
\]
这正是第 \ref{sec:10-3} 节定义~\ref{def:admm-iteration} 的两块线性约束模型，其中 $A=I$、$B=-B$、$b=0$。把定义~\ref{def:admm-iteration} 直接展开到乘积空间，即得到所述更新式。
\end{proof}

\begin{theorem}[分布式推荐模型的收敛模板]\label{thm:distributed-recommendation-convergence}
设定义~\ref{def:distributed-consensus-model} 的问题存在鞍点，且命题~\ref{prop:consensus-admm-splitting} 给出的各步极小化子问题都可解。则相应 ADMM 序列满足：
\begin{enumerate}
  \item 原残差
  \[
  r_i^k:=z_i^k-P_i z^k
  \]
  收敛到 $0$；
  \item 对偶残差收敛到 $0$；
  \item 任意弱聚点都是该共识模型的原--对偶解。
\end{enumerate}
特别地，若解集为单点，则整个序列弱收敛到该解。
\end{theorem}

\begin{proof}
把命题~\ref{prop:consensus-admm-splitting} 中的乘积空间模型视为第 \ref{sec:10-3} 节两块问题的特例。此时
\[
f(x)=\sum_{i=1}^N F_i(z_i),\qquad g(z)=G(z),
\]
而约束算子是有界线性的。于是第 \ref{sec:10-3} 节定理~\ref{thm:admm-primal-dual-convergence} 可直接应用，得到原残差和对偶残差都收敛到零，且每个弱聚点满足相应 KKT 系统。由于这里的 KKT 系统恰是共识约束
\[
z_i=P_i z
\]
与局部/全局次微分最优性条件的组合，故结论成立。
\end{proof}

\begin{corollary}[图推荐与表示匹配的分布式模板]\label{cor:graph-recommendation-distributed-template}
若局部目标取为
\[
F_i(z_i)=\ell_i(z_i)+\tau_i\mathcal E_{\mathscr G_i}(z_i)
\]
或
\[
F_i(z_i)=\mathcal R_{\lambda_i}(M_i),
\]
其中 $\mathcal E_{\mathscr G_i}$ 是第 \ref{sec:12-2} 节定义~\ref{def:graph-dirichlet-energy} 的图能量，$\mathcal R_{\lambda_i}$ 是第 \ref{sec:12-3} 节定义~\ref{def:contrastive-risk-functional} 的对比风险，则相应图推荐、检索匹配和分片 embedding 训练都属于定义~\ref{def:distributed-consensus-model} 的范畴。因而，在定理~\ref{thm:distributed-recommendation-convergence} 的假设下，这些模型可由 ADMM 获得一个统一的分布式收敛模板。
\end{corollary}

\begin{proof}
图能量与对比风险都已在前两节被写成 proper、凸且下半连续泛函；把它们放到局部节点上，再加上共识约束 $z_i=P_i z$，就正好落入定义~\ref{def:distributed-consensus-model}。于是直接应用定理~\ref{thm:distributed-recommendation-convergence} 即可。
\end{proof}

\begin{remark}
从第 \ref{sec:9-4} 节推论~\ref{cor:proximal-point-template} 的角度看，ADMM 仍然是在解一个原--对偶零点问题；只是这里的零点被通信拓扑和参数分片重新组织了。也因此，本节并不扩展到参数服务器、在线排序或工业服务架构细节，而只保留凸分块模型可稳定复用的求解接口。
\end{remark}

\subsection{典型例子}

\begin{example}[共识矩阵分解]
把用户块、物品块或时间块分别放到不同节点上，每个节点维护局部矩阵变量，再通过全局一致性约束保证公共因子对齐。这是分布式矩阵分解最常见的抽象形式。
\end{example}

\begin{example}[图正则推荐]
若每个节点还带有本地图结构，则局部目标中自然出现第 \ref{sec:12-2} 节的图 Dirichlet 能量。这样得到的模型同时耦合了图平滑与全局共识，正是图推荐系统里经常出现的结构。
\end{example}

\begin{example}[分片 embedding 训练]
在大规模推荐系统中，embedding 表往往按词表、用户或物品分片存储。把局部匹配损失写成第 \ref{sec:12-3} 节的对比风险，再通过全局变量协调共享参数，就得到定义~\ref{def:sharded-embedding-problem} 的典型实例。
\end{example}

\subsection*{有限维对照}

Boyd 关于共识 ADMM 的有限维叙述，仍然是理解本节最好的入口：局部子问题、全局平均和残差监控都保留了下来。但本节的关键补充是，推荐系统中的“参数分片”“图正则”“多塔一致性”并不是新的数学对象，它们只是被重新装进了第 \ref{sec:10-3} 节 ADMM 和第 \ref{sec:7-2} 节标准型接口中。换言之，工业推荐系统里的分布式训练并不是脱离理论的工程技巧，而是 Boyd 母语在无穷维与乘积空间中的一次延伸。

\subsection*{习题（待补）}

\begin{exercise}[共识 ADMM 占位]
补一题：从定义~\ref{def:distributed-consensus-model} 出发，完整推出命题~\ref{prop:consensus-admm-splitting} 的三组更新式。
\end{exercise}

\begin{exercise}[图推荐桥接题占位]
补一题：把一个带图平滑正则的推荐模型写成推论~\ref{cor:graph-recommendation-distributed-template} 的形式，并指出局部项与全局一致性项分别来自哪里。
\end{exercise}

